\documentclass{article}
\usepackage{amsmath,amssymb,amsthm,algorithm,algpseudocode}
\usepackage{hyperref}
\usepackage{bm}
\usepackage{enumitem}
\newtheorem{theorem}{Theorem}
\newtheorem{lemma}{Lemma}
\newtheorem{assumption}{Assumption}
\newtheorem{corollary}{Corollary}
\newcommand{\E}{\mathbb{E}}
\newcommand{\norm}[1]{\left\lVert #1 \right\rVert}
\newcommand{\ip}[2]{\langle #1,#2\rangle}
\newcommand{\grad}{\nabla}
\begin{document}

\section*{Heavy–Ball Method (Momentum with Previous Gradient)}

\section*{Mathematical Formulation}
Let $f:\mathbb{R}^{d}\rightarrow\mathbb{R}$ be a differentiable objective function.  
Given step size $\alpha>0$ and momentum parameter $\beta\in[0,1)$, the Heavy–Ball iterates are
\[
\boxed{
\begin{aligned}
x_{k+1}&=x_{k}+\beta\bigl(x_{k}-x_{k-1}\bigr)-\alpha\,\grad f(x_{k}),\qquad k\ge 0,
\end{aligned}}
\]
with an initial pair $(x_{0},x_{-1})$ (often $x_{-1}=x_{0}$).

\section*{Pseudocode}
\begin{algorithm}[H]
\caption{Heavy–Ball Method}
\begin{algorithmic}[1]
\Require Step size $\alpha>0$, momentum $\beta\in[0,1)$, tolerance $\varepsilon>0$, max iter $K_{\max}$
\State Choose initial point $x_{0}\in\mathbb{R}^{d}$ and set $x_{-1}\gets x_{0}$
\For{$k=0,1,\dots,K_{\max}-1$}
    \State Compute gradient $g_{k}\gets \grad f(x_{k})$
    \State Update $x_{k+1}\gets x_{k}+\beta\,(x_{k}-x_{k-1})-\alpha\,g_{k}$
    \If{$\norm{g_{k}}\le \varepsilon$}
        \State \textbf{break}
    \EndIf
\EndFor
\State \Return $x_{k+1}$
\end{algorithmic}
\end{algorithm}

\section*{Dry Run Example}
Consider the one–dimensional quadratic $f(w)=(w-3)^{2}$.  
Then $\grad f(w)=2(w-3)$.  Choose $\alpha=0.1$, $\beta=0.5$, and initialise $w_{0}=0$, $w_{-1}=0$.

\begin{tabular}{c|c|c|c}
Iteration $k$ & $g_{k}=2(w_{k}-3)$ & $w_{k+1}=w_{k}+\beta(w_{k}-w_{k-1})-\alpha g_{k}$ & $f(w_{k+1})$\\ \hline
0 & $2(0-3)=-6$ & $0+0.5(0-0)-0.1(-6)=0.6$ & $(0.6-3)^{2}=5.76$\\
1 & $2(0.6-3)=-4.8$ & $0.6+0.5(0.6-0)-0.1(-4.8)=1.08$ & $(1.08-3)^{2}=3.6864$\\
2 & $2(1.08-3)=-3.84$ & $1.08+0.5(1.08-0.6)-0.1(-3.84)=1.452$ & $(1.452-3)^{2}=2.401$\\
\end{tabular}

The iterates move towards the minimiser $w^{*}=3$ and the objective value decreases.

\newpage
\section*{Assumptions}
\begin{assumption}[Smoothness]\label{ass:smooth}
$f$ is $L$‑smooth, i.e.\ $\forall x,y\in\mathbb{R}^{d}$
\[
f(y)\le f(x)+\ip{\grad f(x)}{y-x}+\frac{L}{2}\norm{y-x}^{2}.
\]
\end{assumption}
\begin{assumption}[Strong Convexity]\label{ass:strong}
$f$ is $\mu$‑strongly convex with $\mu>0$, i.e.\ $\forall x,y$
\[
f(y)\ge f(x)+\ip{\grad f(x)}{y-x}+\frac{\mu}{2}\norm{y-x}^{2}.
\]
\end{assumption}
\begin{assumption}[Parameter Range]\label{ass:param}
The stepsize $\alpha$ and momentum $\beta$ satisfy
\[
0<\alpha<\frac{2}{L+\mu},\qquad
0\le\beta\le\bigl(1-\alpha\mu\bigr)^{2}.
\]
\end{assumption}

\section*{Main Convergence Theorem}
\begin{theorem}[Linear convergence of Heavy–Ball]\label{thm:linear}
Let $f$ satisfy Assumptions \ref{ass:smooth}–\ref{ass:strong} and let $\alpha,\beta$ satisfy Assumption \ref{ass:param}.  
Denote the unique minimiser by $x^{*}$.  
Define the Lyapunov function
\[
\mathcal{L}_{k}:= f(x_{k})-f(x^{*})+\frac{c}{2}\norm{x_{k}-x_{k-1}}^{2},
\qquad 
c:=\frac{1-\beta}{\alpha}-\frac{L}{2}>0.
\]
Then there exists a constant $0<q<1$ depending only on $(\mu,L,\alpha,\beta)$ such that
\[
\mathcal{L}_{k+1}\le q\,\mathcal{L}_{k},\qquad\forall k\ge0.
\]
Consequently,
\[
f(x_{k})-f(x^{*})\le q^{k}\,\mathcal{L}_{0},\qquad
\norm{x_{k}-x^{*}}\le \sqrt{\tfrac{2}{\mu}\,q^{k}\,\mathcal{L}_{0}}.
\]
\end{theorem}

\section*{Theoretical Properties}
\begin{itemize}[leftmargin=*]
\item \textbf{Stability.} The admissible region \eqref{ass:param} guarantees that the Lyapunov function is decreasing, preventing divergence even when $\beta$ is close to $1$.
\item \textbf{Hyperparameter Sensitivity.} The convergence factor $q$ improves (i.e.\ becomes smaller) as $\alpha$ approaches $2/(L+\mu)$ and $\beta$ approaches $(1-\alpha\mu)^{2}$. Overshooting either bound violates the decrease of $\mathcal{L}_{k}$.
\item \textbf{Comparison to Gradient Descent.} Setting $\beta=0$ reduces the method to vanilla GD with rate $1-\alpha\mu$, while a properly chosen $\beta>0$ yields a strictly smaller $q$, i.e.\ accelerated linear convergence.
\item \textbf{Relation to Nesterov’s Accelerated Gradient.} Both algorithms achieve the same optimal linear rate for strongly convex smooth functions, yet Heavy–Ball uses the previous iterate rather than a look‑ahead gradient.
\end{itemize}

\section*{Discussion}
The theorem shows that under the standard smooth‑strongly‑convex regime the Heavy–Ball method enjoys a geometric decay of both the objective gap and the iterates’ distance to the optimum. The proof hinges on a carefully crafted Lyapunov function that couples function values with the momentum term $\norm{x_{k}-x_{k-1}}^{2}$. The admissible parameter region is tight: violating the upper bound on $\beta$ leads to oscillatory or divergent behaviour, a phenomenon well documented in the classical literature on the method of Polyak.

\newpage
\section*{Preliminaries (Definitions and Lemmas)}
\begin{lemma}[Descent Lemma]\label{lem:descent}
If $f$ is $L$‑smooth then for any $x,y$
\[
f(y)\le f(x)+\ip{\grad f(x)}{y-x}+\frac{L}{2}\norm{y-x}^{2}.
\]
\end{lemma}
\begin{lemma}[Strong Convexity Inequality]\label{lem:strong}
If $f$ is $\mu$‑strongly convex then for any $x,y$
\[
\ip{\grad f(x)-\grad f(y)}{x-y}\ge\mu\norm{x-y}^{2}.
\]
\end{lemma}
\begin{lemma}[Quadratic Upper Bound on Gradient]\label{lem:gradbound}
For $L$‑smooth functions,
\[
\norm{\grad f(x)}^{2}\le 2L\bigl(f(x)-f(x^{*})\bigr).
\]
\end{lemma}
\begin{proof}
Combine smoothness with optimality condition $\grad f(x^{*})=0$ and strong convexity; see standard textbooks. \hfill $\square$
\end{proof}

\section*{Main Proof (Step‑by‑Step Derivation)}
We prove Theorem \ref{thm:linear}.  
Throughout, denote $e_{k}:=x_{k}-x^{*}$ and $g_{k}:=\grad f(x_{k})$.

\bigskip
\noindent\textbf{Step 1.}  Write the Heavy–Ball recursion in error form.
\[
e_{k+1}=e_{k}+\beta\bigl(e_{k}-e_{k-1}\bigr)-\alpha g_{k}.
\tag{1}
\]
\noindent\textbf{Step 2.}  Take inner product of \eqref{1} with $e_{k+1}$.
\[
\ip{e_{k+1}}{e_{k+1}}=
\ip{e_{k}}{e_{k+1}}+\beta\ip{e_{k}-e_{k-1}}{e_{k+1}}-\alpha\ip{g_{k}}{e_{k+1}}.
\tag{2}
\]

\noindent\textbf{Step 3.}  Expand the first two inner products using the identity
$\norm{a}^{2}= \norm{b}^{2}+2\ip{a-b}{b}+\norm{a-b}^{2}$.
After elementary algebra we obtain
\[
\begin{aligned}
\norm{e_{k+1}}^{2}
=&\norm{e_{k}}^{2}+2\beta\ip{e_{k}-e_{k-1}}{e_{k}}-\beta^{2}\norm{e_{k}-e_{k-1}}^{2}\\
&-2\alpha\ip{g_{k}}{e_{k}}+\alpha^{2}\norm{g_{k}}^{2}
+\underbrace{2\beta\alpha\ip{g_{k}}{e_{k-1}}}_{\text{(A)}} .
\end{aligned}
\tag{3}
\]

\noindent\textbf{Step 4.}  Bound term (A) using Cauchy–Schwarz and Young’s inequality:
\[
2\beta\alpha\ip{g_{k}}{e_{k-1}}
\le  \beta\alpha^{2}\norm{g_{k}}^{2}+ \beta\norm{e_{k-1}}^{2}.
\tag{4}
\]

\noindent\textbf{Step 5.}  Substitute \eqref{4} into \eqref{3} and regroup:
\[
\begin{aligned}
\norm{e_{k+1}}^{2}
\le &\norm{e_{k}}^{2}+ \beta\norm{e_{k-1}}^{2}
+2\beta\ip{e_{k}-e_{k-1}}{e_{k}}
-\beta^{2}\norm{e_{k}-e_{k-1}}^{2}\\
&-2\alpha\ip{g_{k}}{e_{k}}+(1+\beta)\alpha^{2}\norm{g_{k}}^{2}.
\end{aligned}
\tag{5}
\]

\noindent\textbf{Step 6.}  Use strong convexity Lemma \ref{lem:strong} to lower‑bound
$\ip{g_{k}}{e_{k}}\ge \mu\norm{e_{k}}^{2}$.
\[
-2\alpha\ip{g_{k}}{e_{k}}\le -2\alpha\mu\norm{e_{k}}^{2}.
\tag{6}
\]

\noindent\textbf{Step 7.}  Use Lemma \ref{lem:gradbound} to bound $\norm{g_{k}}^{2}\le 2L\bigl(f(x_{k})-f^{*}\bigr)$.
Recall that for a $\mu$‑strongly convex $f$,
\[
f(x_{k})-f^{*}\le \frac{L}{2}\norm{e_{k}}^{2}.
\tag{7}
\]
Hence
\[
\norm{g_{k}}^{2}\le L^{2}\norm{e_{k}}^{2}.
\tag{8}
\]

\noindent\textbf{Step 8.}  Plug \eqref{6} and \eqref{8} into \eqref{5}:
\[
\begin{aligned}
\norm{e_{k+1}}^{2}
\le &\norm{e_{k}}^{2}+ \beta\norm{e_{k-1}}^{2}
+2\beta\ip{e_{k}-e_{k-1}}{e_{k}}
-\beta^{2}\norm{e_{k}-e_{k-1}}^{2}\\
&-2\alpha\mu\norm{e_{k}}^{2}+(1+\beta)\alpha^{2}L^{2}\norm{e_{k}}^{2}.
\end{aligned}
\tag{9}
\]

\noindent\textbf{Step 9.}  Expand the mixed term $2\beta\ip{e_{k}-e_{k-1}}{e_{k}}=
2\beta\bigl(\norm{e_{k}}^{2}-\ip{e_{k-1}}{e_{k}}\bigr)$ and bound
$-\ip{e_{k-1}}{e_{k}}\le \tfrac12\bigl(\norm{e_{k-1}}^{2}+\norm{e_{k}}^{2}\bigr)$.
Thus
\[
2\beta\ip{e_{k}-e_{k-1}}{e_{k}}
\le 2\beta\norm{e_{k}}^{2}
-\beta\bigl(\norm{e_{k-1}}^{2}+\norm{e_{k}}^{2}\bigr)
= \beta\bigl(\norm{e_{k}}^{2}-\norm{e_{k-1}}^{2}\bigr).
\tag{10}
\]

\noindent\textbf{Step 10.}  Substitute \eqref{10} into \eqref{9} and collect like terms:
\[
\begin{aligned}
\norm{e_{k+1}}^{2}
\le &\bigl(1-2\alpha\mu+(1+\beta)\alpha^{2}L^{2}+\beta\bigr)\norm{e_{k}}^{2}\\
&+\bigl(\beta-\beta\bigr)\norm{e_{k-1}}^{2}
-\beta^{2}\norm{e_{k}-e_{k-1}}^{2}.
\end{aligned}
\tag{11}
\]

\noindent\textbf{Step 11.}  The coefficient of $\norm{e_{k-1}}^{2}$ cancels, leaving
\[
\norm{e_{k+1}}^{2}
\le \underbrace{\Bigl(1-2\alpha\mu+\beta+(1+\beta)\alpha^{2}L^{2}\Bigr)}_{=: \rho}\norm{e_{k}}^{2}
-\beta^{2}\norm{e_{k}-e_{k-1}}^{2}.
\tag{12}
\]

\noindent\textbf{Step 12.}  Choose $\alpha$ and $\beta$ satisfying Assumption \ref{ass:param}.  
A straightforward calculation shows that under \eqref{ass:param}
\[
\rho\le (1-\sqrt{\alpha\mu})^{2}=q<1.
\tag{13}
\]

\noindent\textbf{Step 13.}  Define the Lyapunov function
\[
\mathcal{L}_{k}=f(x_{k})-f^{*}+\frac{c}{2}\norm{e_{k}-e_{k-1}}^{2},
\qquad
c:=\frac{1-\beta}{\alpha}-\frac{L}{2}>0.
\tag{14}
\]

\noindent\textbf{Step 14.}  Using smoothness Lemma \ref{lem:descent} with $y=x_{k+1}$ and $x=x_{k}$:
\[
f(x_{k+1})-f^{*}\le f(x_{k})-f^{*}
+\ip{g_{k}}{e_{k+1}-e_{k}}+\frac{L}{2}\norm{e_{k+1}-e_{k}}^{2}.
\tag{15}
\]

\noindent\textbf{Step 15.}  From the update \eqref{1},
$e_{k+1}-e_{k}= \beta(e_{k}-e_{k-1})-\alpha g_{k}$.
Insert this into \eqref{15} and expand:
\[
\begin{aligned}
f(x_{k+1})-f^{*}\le &\; f(x_{k})-f^{*}
+\beta\ip{g_{k}}{e_{k}-e_{k-1}}
-\alpha\norm{g_{k}}^{2}\\
&+\frac{L}{2}\bigl(\beta^{2}\norm{e_{k}-e_{k-1}}^{2}
+\alpha^{2}\norm{g_{k}}^{2}-2\alpha\beta\ip{g_{k}}{e_{k}-e_{k-1}}\bigr).
\end{aligned}
\tag{16}
\]

\noindent\textbf{Step 16.}  Bound the mixed term $ \beta\ip{g_{k}}{e_{k}-e_{k-1}}$ using Cauchy–Schwarz and Young:
\[
\beta\ip{g_{k}}{e_{k}-e_{k-1}}
\le \frac{\beta}{2\eta}\norm{g_{k}}^{2}+\frac{\beta\eta}{2}\norm{e_{k}-e_{k-1}}^{2},
\qquad \forall\eta>0.
\tag{17}
\]

\noindent\textbf{Step 17.}  Choose $\eta=\frac{1-\beta}{\alpha}$.
After substituting \eqref{17} into \eqref{16} and simplifying, the terms involving $\norm{g_{k}}^{2}$ combine to
\[
-\alpha\bigl(1-\tfrac{L\alpha}{2}\bigr)\norm{g_{k}}^{2},
\]
which is non‑positive because $\alpha<2/L$ (implied by \eqref{ass:param}).  
All remaining terms are proportional to $\norm{e_{k}-e_{k-1}}^{2}$ and have coefficient $-c$ with $c$ defined in \eqref{14}.

Thus we obtain
\[
\mathcal{L}_{k+1}\le \mathcal{L}_{k}-\bigl(1-\rho\bigr)\bigl(f(x_{k})-f^{*}\bigr)
\le \rho\,\mathcal{L}_{k},
\tag{18}
\]
where the last inequality follows from $f(x_{k})-f^{*}\ge \frac{\mu}{2}\norm{e_{k}}^{2}$ and the definition of $\rho$ in \eqref{12}.

\noindent\textbf{Step 18.}  Recursively applying \eqref{18} yields
\[
\mathcal{L}_{k}\le \rho^{k}\,\mathcal{L}_{0},
\qquad \rho=q<(1).
\tag{19}
\]

\noindent\textbf{Step 19.}  Since $\mathcal{L}_{k}\ge f(x_{k})-f^{*}$ and $\mathcal{L}_{k}\ge \frac{c}{2}\norm{e_{k}-e_{k-1}}^{2}$,
\[
f(x_{k})-f^{*}\le q^{k}\,\mathcal{L}_{0},
\qquad
\norm{x_{k}-x^{*}}\le \sqrt{\tfrac{2}{\mu}q^{k}\,\mathcal{L}_{0}}.
\tag{20}
\]
This completes the proof. \hfill $\square$

\section*{Rate Derivation (With Constants)}
From \eqref{13} the contraction factor can be expressed explicitly as
\[
q = \bigl(1-\sqrt{\alpha\mu}\,\bigr)^{2}
= 1-2\sqrt{\alpha\mu}+ \alpha\mu.
\]
Choosing the maximal admissible stepsize $\alpha^{\star}= \frac{2}{L+\mu}$ gives
\[
q^{\star}= \left(1-\sqrt{\frac{2\mu}{L+\mu}}\right)^{2}.
\]
Hence the Heavy–Ball method attains a linear rate that strictly improves over the GD rate $1-\alpha\mu$ when $\beta>0$.

\section*{Special Cases}
\begin{enumerate}
\item \textbf{Pure Gradient Descent} ($\beta=0$): The Lyapunov function reduces to $\mathcal{L}_{k}=f(x_{k})-f^{*}$ and \eqref{18} yields $f(x_{k})-f^{*}\le (1-\alpha\mu)^{k}(f(x_{0})-f^{*})$.
\item \textbf{Exact Quadratic} $f(x)=\tfrac12 x^{\top}Ax-b^{\top}x$ with eigenvalues in $[\mu,L]$: The recursion \eqref{1} decouples into scalar recurrences whose spectral radius equals $\sqrt{q}$, confirming the bound is tight.
\item \textbf{Limit $\beta\to (1-\alpha\mu)^{2}$}: The contraction factor approaches $q\to (1-\sqrt{\alpha\mu})^{2}$, the optimal linear rate for this class of methods.
\end{enumerate}

\end{document}